% Part: normal-modal-logic
% Chapter: syntax-and-semantics
% Section: language-modal-logic

\documentclass[../../../include/open-logic-section]{subfiles}

\begin{document}

\olfileid{nml}{syn}{lan}

\olsection{لغة المنطق الموجهي الأساسي}

\begin{defn}
تحتوي اللغة الأساسية للمنطق الموجهي على
\begin{enumerate}
  \tagitem{prvFalse}{ثابت القضايا للدلالة على !!{falsity}~$\lfalse$.}{}
  \tagitem{prvTrue}{ثابت القضايا للدلالة على !!{truth}~$\ltrue$.}{}
  \item مجموعة !!{denumerable}s من !!{propositional variable}s:
    $\Obj p_0$، و$\Obj p_1$، و$\Obj p_2$، و\dots
  \item روابط القضايا: \startycommalist
  \iftag{prvNot}{\ycomma $\lnot$ (النفي)}{}%
  \iftag{prvAnd}{\ycomma $\land$ (الاقتران)}{}%
  \iftag{prvOr}{\ycomma $\lor$ (الفصل)}{}%
  \iftag{prvIf}{\ycomma $\lif$ (!!{conditional})}{}%
  \iftag{prvIff}{\ycomma $\liff$ (!!{biconditional})}{}.
  \tagitem{prvBox}{المؤثر الموجهي~$\Box$.}{}
  \tagitem{prvDiamond}{المؤثر الموجهي~$\Diamond$.}{}
\end{enumerate}
\end{defn}

\begin{defn}
تُعرّف \emph{!!^{formula}s} اللغة الموجهية الأساسية استقرائيًا كما
يأتي:
\begin{enumerate}
\tagitem{prvFalse}{$\lfalse$ !!{formula} ذرية.}{}

\tagitem{prvTrue}{$\ltrue$ !!{formula} ذرية.}{}

\item كل !!{propositional variable}~$\Obj p_i$ !!{formula} (ذرية).

\tagitem{prvNot}{إذا كانت~$!A$ !!a{formula}، فإن~$\lnot !A$
  !!a{formula}.}{}

\tagitem{prvAnd}{إذا كانت~$!A$ و$!B$ !!{formula}s، فإن
  $(!A \land !B)$ !!a{formula}.}{}

\tagitem{prvOr}{إذا كانت~$!A$ و$!B$ !!{formula}s، فإن
  $(!A \lor !B)$ !!a{formula}.}{}

\tagitem{prvIf}{إذا كانت~$!A$ و$!B$ !!{formula}s، فإن
  $(!A \lif !B)$ !!a{formula}.}{}

\tagitem{prvIff}{إذا كانت~$!A$ و$!B$ !!{formula}s، فإن
  $(!A \liff !B)$ !!a{formula}.}{}

\tagitem{prvBox}{إذا كانت~$!A$ !!a{formula}، فإن~$\Box !A$
  !!a{formula}.}{}

\tagitem{prvDiamond}{إذا كانت~$!A$ !!a{formula}، فإن~$\Diamond !A$
  !!a{formula}.}{}

\tagitem{limitClause}{لا شيء آخر !!a{formula}.}{}
\end{enumerate}
\end{defn}

\begin{tagblock}{defNot,defOr,defAnd,defIf,defIff,defTrue,defFalse,defBox,defDiamond}
\begin{defn}
تُفهم الصيغ المبنية بالمؤثرات المعرّفة كما يأتي:

\begin{tagenumerate}{defTrue,defFalse,defNot,defOr,defAnd,defIf,defIff,defBox,defDiamond}

\tagitem{defTrue}{$\ltrue$ اختصار لـ
  \iftag{prvFalse}{$\lnot\lfalse$}{$(!A \lor \lnot !A)$ لبعض
    !!{formula} ذرية ثابتة~$!A$}.}{}

\tagitem{defFalse}{$\lfalse$ اختصار لـ
  \iftag{prvTrue}{$\lnot\ltrue$}{$(!A \land \lnot !A)$ لبعض
    !!{formula} ذرية ثابتة~$!A$}.}{}

\tagitem{defNot}{$\lnot !A$ اختصار لـ$!A \lif \lfalse$.}{}

\tagitem{defOr}{$!A \lor !B$ اختصار لـ
  \iftag{prvAnd}{$\lnot(\lnot !A \land \lnot !B)$}{$\lnot !A \lif
    !B$}.}{}

\tagitem{defAnd}{$!A \land !B$ اختصار لـ
  \iftag{prvOr}{$\lnot(\lnot !A \lor \lnot !B)$}{$\lnot (!A \lif
    \lnot !B)$}.}{}

\tagitem{defIf}{$!A \lif !B$ اختصار لـ
  \iftag{prvOr}{$(\lnot !A \lor !B)$}{$\lnot (!A \land \lnot !B)$}.}{}

\tagitem{defIff}{$!A \liff !B$ اختصار لـ$(!A \lif !B) \land (!B
  \lif !A)$.}{}

\tagitem{defBox}{$\Box !A$ اختصار لـ$\lnot\Diamond\lnot !A$. }{}

\tagitem{defDiamond}{$\Diamond !A$ اختصار لـ$\lnot\Box\lnot !A$.}{}
\end{tagenumerate}
\end{defn}
\end{tagblock}

إذا كانت !!a{formula}~$!A$ لا تحتوي~$\Box$ ولا~$\Diamond$، قلنا إنها
\emph{خالية من الجهات}.

\end{document}
